\begin{statementbox}
\par\medskip
\noindent\begin{minipage}{\linewidth}
\paragraph{F081 罗尔定理与零点个数关系}
\begin{center}% F081: equal endpoint heights and an interior horizontal tangent.
\begin{tikzpicture}[x=1cm,y=1cm,>=stealth,font=\small]
\draw[rounded corners=3pt,gray!40] (0,0) rectangle (14.2,5.5);
\node at (7.1,4.9) {两端等高，内部至少有一条水平切线};
\begin{scope}[shift={(3.7,1.5)},x=2cm,y=1.7cm]
\draw[->,gray!65] (-1.4,0)--(1.5,0) node[right] {$x$};
\draw[->,gray!65] (0,-.25)--(0,1.45) node[above] {$y$};
\draw[blue!75!black,line width=1.15pt,domain=-1:1,samples=120] plot (\x,{1-(\x)*(\x)});
\fill[blue!75!black] (-1,0) circle (1.8pt) (1,0) circle (1.8pt);
\node[below] at (-1,0) {$A(a,f(a))$};
\node[below] at (1,0) {$B(b,f(b))$};
\draw[orange!85!black,line width=1.2pt] (-.35,1)--(.35,1);
\fill[orange!85!black] (0,1) circle (1.8pt);
\node[above left] at (0,1) {$C$};
\draw[dashed,gray!60] (0,0)--(0,1);\node[below right] at (0,0) {$c$};
\end{scope}
\node at (10.5,3.6) {$f(a)=f(b)$};
\node at (10.5,2.6) {$a<c<b$};
\node[orange!85!black] at (10.5,1.6) {$f'(c)=0$};
\node at (7.1,.5) {前提：在 $[a,b]$ 连续，在 $(a,b)$ 可导。};
\end{tikzpicture}
\end{center}
\end{minipage}\par\medskip
\textbf{一句话先记住：两头一样高，中间至少有一处走平。}

设 $a<b$，如果函数 $f(x)$ 满足：
\begin{enumerate}
  \item 从 $a$ 到 $b$ 的图像\textbf{连在一起，没有断开}（在 $[a,b]$ 上连续）；
  \item 在 $a$ 与 $b$ 之间，每一点\textbf{都能求导}（在 $(a,b)$ 内可导）；
  \item 两端的函数值相等，即 $f(a)=f(b)$；
\end{enumerate}
那么一定能在两端之间找到至少一个 $c$，使
\[
\boxed{\ a<c<b,\qquad f'(c)=0\ }.
\]

\textbf{用图形语言说：}令 $h=f(a)=f(b)$，取
\[
A(a,h),\qquad B(b,h),\qquad C(c,f(c)).
\]
因为 $A,B$ 同高，直线 $AB$ 是水平的，其斜率为
\[
 k_{AB}=\frac{f(b)-f(a)}{b-a}=0.
\]
而 $f'(c)=0$ 的几何意思是：\textbf{点 $C$ 处的切线也水平}，所以
\[
\boxed{\ \text{点 }C\text{ 处的切线}\ \parallel\ AB\ \parallel\ x\text{ 轴}\ }.
\]
\textbf{注意：}切线与 $AB$ 平行，不代表切点 $C$ 一定和 $A,B$ 一样高。

这就是\textbf{罗尔定理}。

\textbf{它和零点有什么关系？}如果 $f(a)=f(b)=0$，那么曲线在 $a,b$ 两处碰到或穿过 $x$ 轴，中间必有一处切线水平。
\end{statementbox}
\clearpage
\begin{statementbox}
\paragraph{真正实用的结论：零点越多，导数为零的位置越多}
设 $f$ 在一个区间 $I$ 上连续，在 $I$ 的内部可导，并且有 $m\ge2$ 个\textbf{互不相同}的零点：
\[
x_1<x_2<\cdots<x_m,\qquad f(x_1)=\cdots=f(x_m)=0.
\]
\begin{center}% F081: precise mathematical teaching diagram.
\begin{tikzpicture}[x=1cm,y=1cm,>=stealth,font=\small]
\draw[rounded corners=3pt,gray!40] (0,0) rectangle (14.2,9.2);
\begin{scope}[yshift=1cm]
\node at (7.1,7.6) {四个不同零点，夹出三个不同的导数零点};
\begin{scope}[shift={(7.1,4.2)},x=1.65cm,y=0.16cm]
\draw[->,gray!65] (-3.6,0)--(3.6,0) node[right] {$x$};
\draw[->,gray!65] (0,-19)--(0,14) node[above] {$y$};
\draw[blue!75!black,line width=1.15pt,,domain=-3.12:3.12,samples=160] plot (\x,{(\x)*(\x)*(\x)*(\x)-10*(\x)*(\x)+9});
\fill[blue!75!black] (-3,0) circle (1.8pt);\node[above right] at (-3,0) {$x_1$};
\fill[blue!75!black] (-1,0) circle (1.8pt);\node[above left] at (-1,0) {$x_2$};
\fill[blue!75!black] (1,0) circle (1.8pt);\node[above right] at (1,0) {$x_3$};
\fill[blue!75!black] (3,0) circle (1.8pt);\node[above left] at (3,0) {$x_4$};
\draw[orange!85!black,line width=1.2pt] (-2.5560679774997896,-16.000000000000004)--(-1.9160679774997897,-16.000000000000004);
\fill[orange!85!black] (-2.23606797749979,-16.000000000000004) circle (1.8pt);\node[below] at (-2.23606797749979,-16.000000000000004) {$c_1$};
\draw[orange!85!black,line width=1.2pt] (-0.32,9)--(0.32,9);
\fill[orange!85!black] (0,9) circle (1.8pt);\node[above left] at (0,9) {$c_2$};
\draw[orange!85!black,line width=1.2pt] (1.9160679774997897,-16.000000000000004)--(2.5560679774997896,-16.000000000000004);
\fill[orange!85!black] (2.23606797749979,-16.000000000000004) circle (1.8pt);\node[below] at (2.23606797749979,-16.000000000000004) {$c_3$};
\end{scope}
\node[] at (7.1,.65) {$x_1<c_1<x_2<c_2<x_3<c_3<x_4$};
\end{scope}
\node[text width=13cm,align=center] at (7.1,.9) {图例 $f(x)=(x^2-1)(x^2-9)$；每个开区间各给一个，位置不会重复。};
\end{tikzpicture}
\end{center}
在每一对相邻零点之间，至少有一个导数为 $0$ 的位置。
这些位置分处不同的小区间，因此\textbf{互不相同}。于是
\[
\boxed{\ f(x)\text{ 有 }m\text{ 个不同零点}\ \Longrightarrow\ f'(x)\text{ 至少有 }m-1\text{ 个不同零点}\ }.
\]
\textbf{反过来用往往更快：}设 $k$ 为非负整数，则
\[
\boxed{\ f'(x)\text{ 最多有 }k\text{ 个零点}\ \Longrightarrow\ f(x)\text{ 最多有 }k+1\text{ 个零点}\ }.
\]
特别地，若 $f'(x)$ 在该区间内\textbf{没有零点}，则 $f(x)=0$ 在该区间内\textbf{最多只有一个解}。

\textbf{注意：}“最多”并不等于“一定有”。想证明\textbf{恰好}几个零点，还必须找到确实存在这些零点的理由。

\textbf{眼前一亮的推广：不只是数与 $x$ 轴的交点。}把 $x$ 轴换成\textbf{任意水平直线 $y=h$}，结论仍然成立。如果图像与 $y=h$ 有 $m$ 个不同交点，中间至少有 $m-1$ 处水平切线；如果整个区间最多只有 $k$ 个水平切线位置，那么\textbf{任意一条水平直线最多与图像相交 $k+1$ 次}。这里的数量都按不同的横坐标计算。

\end{statementbox}
